\section{Threats to Validity}
\label{sec:threats}

\paragraph{IR count as the objective.}
All four hypotheses use IR instruction count. We confirm the H1 headroom in \texttt{.text} size on
two ISAs. For executed instructions and wall-clock time, the evidence is only suggestive. Searched
sequences that minimise IR count are slower than \Othree{} (\S\ref{sec:discussion}).

\paragraph{One compiler version, one machine.}
All results come from LLVM~\LLVMVersion{} on one x86-64 laptop. AArch64 numbers are compile-only
code size. Other LLVM versions change the pass pipeline and could change the interactions.

\paragraph{Timing noise.}
Timing ran under WSL2 on a laptop: serial, pinned to one core, with 20 runs per binary in a
randomised order. Across \NRuntimeBinaries{} timed binaries the median coefficient of variation is
\RuntimeCV\%, but \RuntimeNoisy{} binaries (\RuntimeNoisyPct\%) exceed 5\%. Background load on the
host varied during the timing stage. Randomised order spreads that drift across methods, but it
adds noise~\cite{mytkowicz2009}. No hypothesis depends on timing.

\paragraph{Reimplemented baselines.}
DCO-style and Pan-style are our reimplementations of the published ideas under a common budget and
evaluator, not the original systems. We do not compare against the numbers in those papers.

\paragraph{Corpus.}
\NPrograms{} programs from four suites, \NOwn{} of them small kernels. Results may differ on
large applications. Leave-one-program-out evaluation keeps each program unseen by its own model and
graphs. Programs from the same suite can still be similar, so the static graph shares that
advantage.
